% Target: rmp-iii-a-spt-mpo
% Formula: Pulling O_g through A deposits the physical action U_g on A.
% Ink: Canonical public kernel PEPS star, crossing boxes, and rotated arcs.
% Boundary: Four virtual legs, one physical leg, and two g-string ends are open.
% Source: paper TN-Review-main.tex line 1459; author source ImagesReview Section
%   III/ImagesReview Section III.tex lines 414-444
% Capabilities: kernel, pulling-through, crossing-order, rotated-action
\[
  \tenkzkernel
  \tndeclare{species}{lattice}{hue=source:black}
  \tndeclare{species}{g}{hue=source:red}
  \begin{tenkzeq}[size=m]
    \begin{tenkz}[rows={wire,wire,wire}, cols=3, bonds=none]
      % Four virtual legs end at A; its typed physical port stays source-short.
      \tn[at=(2,2), name=A,
        ports={0:virtual,45:virtual,90:physical,180:virtual,225:virtual}]{}
      \tnwire[kind=index, species=lattice, name=vxw]{(2,1)}{A.180}
      \tnwire[kind=index, species=lattice, name=vxe]{A.0}{(2,3)}
      \tnwire[kind=index, species=lattice, name=vdsw]{(3,1)}{A.225}
      \tnwire[kind=index, species=lattice, name=vdne]{A.45}{(1,3)}
      % the g string wraps the northwest corner, crossing two legs
      \tn[at=midway (1,1) and (2,2), skin=none, name=jg1]{}
      \tn[at=midway (2,2) and (1,3), skin=none, name=jg2]{}
      \tnwire[kind=string, species=g, name=g1,
              cross={over at crossing of g1 and vxw}]
        {midway (2,1) and (3,1)}{jg1}
      \tnwire[kind=string, species=g, name=g2,
              cross={over at crossing of g2 and vdne,
                     over at crossing of g1 and g2}]
        {jg1}{jg2}
      \tnwire[kind=string, species=g, name=g3,
              cross={over at crossing of g3 and vdne}]
        {jg2}{midway (1,3) and (2,3)}
      \tn[skin=box, at=crossing of g1 and vxw]{}
      \tn[skin=box, at=crossing of g2 and vdne]{}
      \tnmark[form=label, label pos=nw]{on g1 0.8195}{$g$}
    \end{tenkz}
    =
    \begin{tenkz}[rows={wire,wire,wire}, cols=3, bonds=none]
      \tn[at=(2,2), name=A,
        ports={0:virtual,45:virtual,90:physical,180:virtual,225:virtual}]{}
      \tnwire[kind=index, species=lattice, name=vxw]{(2,1)}{A.180}
      \tnwire[kind=index, species=lattice, name=vxe]{A.0}{(2,3)}
      \tnwire[kind=index, species=lattice, name=vdsw]{(3,1)}{A.225}
      \tnwire[kind=index, species=lattice, name=vdne]{A.45}{(1,3)}
      \tn[at=(1,2), name=Ug, ports={270:physical,90:physical}, label pos=w]{U_g}
      \tnwire[name=wp]{A.90}{Ug.270}
      % the same string, pulled through to the southeast corner
      \tn[at=midway (2,2) and (3,3), skin=none, name=jg1]{}
      \tn[at=midway (2,2) and (3,1), skin=none, name=jg2]{}
      \tnwire[kind=string, species=g, name=g1,
              cross={over at crossing of g1 and vxe}]
        {midway (1,3) and (2,3)}{jg1}
      \tnwire[kind=string, species=g, name=g2,
              cross={over at crossing of g2 and vdsw,
                     over at crossing of g1 and g2}]
        {jg1}{jg2}
      \tnwire[kind=string, species=g, name=g3,
              cross={over at crossing of g3 and vdsw}]
        {jg2}{midway (2,1) and (3,1)}
      \tn[skin=box, at=crossing of g1 and vxe]{}
      \tn[skin=box, at=crossing of g2 and vdsw]{}
      \tnmark[form=label, label pos=se]{on g2 0.5837}{$g$}
    \end{tenkz}
  \end{tenkzeq}.
\]
